\documentclass[aspectratio=169]{beamer}

% \usetheme{CambridgeUS}
% \usecolortheme{default}

% PREVIOUS RED THEME
% \usetheme{Madrid}
% \usecolortheme{beaver}

\setbeamertemplate{frametitle}[default][center] % Centering of slide heading

%%%%%%%%% CUSTOM THEME %%%%%%%%%%
% Found in: https://ramblingacademic.com/2015/12/08/how-to-quickly-overhaul-beamer-colors/

\usetheme{Madrid}
% \useoutertheme{miniframes} % Alternatively: miniframes, infolines, split
\useinnertheme{circles}

% ORIG COLORS
% \definecolor{UBCblue}{rgb}{0.04706, 0.13725, 0.26667} % UBC Blue (primary)
% \definecolor{UBCgrey}{rgb}{0.3686, 0.5255, 0.6235} % UBC Grey (secondary)

\definecolor{UBCblue}{rgb}{0.4823, 0.0, 0.2471} % VU (primary)
\definecolor{UBCgrey}{rgb}{0.8627, 0.8627, 0.8627} % VU (secondary)

%R:123 G:0 B:63

\setbeamercolor{palette primary}{bg=UBCblue,fg=white}
\setbeamercolor{palette secondary}{bg=UBCblue,fg=white}
\setbeamercolor{palette tertiary}{bg=UBCblue,fg=white}
\setbeamercolor{palette quaternary}{bg=UBCblue,fg=white}
\setbeamercolor{structure}{fg=UBCblue} % itemize, enumerate, etc
\setbeamercolor{section in toc}{fg=UBCblue} % TOC sections

% Override palette coloring with secondary
\setbeamercolor{subsection in head/foot}{bg=UBCgrey,fg=white}

% TODO LIST for conference
% 1. unify colours
% 2. use MIF colour palette or VU colour palette
% 3. ...

% \addtobeamertemplate{navigation symbols}{}{%
%     \usebeamerfont{footline}%
%     \usebeamercolor[fg]{footline}%
%     \hspace{1em}%
%     \insertframenumber/\inserttotalframenumber
% }

\usepackage[utf8x]{inputenc}
\def\LTfontencoding{L7x}
\PrerenderUnicode{ąčęėįšųūž}
\usepackage{times}
\usepackage[T1]{fontenc}
\usepackage{color}
\usepackage{verbatim}
\usepackage{graphicx}
\usepackage{fancyvrb}
\usepackage{bm}
\usepackage{amsfonts}
\usepackage{float}
\usepackage{hyperref}
\usepackage{tikz}
\usepackage{textpos} % FOR LOGO
% \usepackage[english]{babel} # THE SLIDES ARE IN LITHUANIAN
\usepackage{LTPlius} % COMMENT IF ENGLISH
\usepackage{tabularx,booktabs}
% \usepackage{enumitem}% http://ctan.org/pkg/enumitem

% set numbers to image, tables, ... captions.
% From: https://tex.stackexchange.com/a/216437
\setbeamertemplate{caption}[numbered]{}% Number float-like environments

\usepackage{caption} % <<< does not include "figure ..."
% \usepackage[labelformat=empty]{caption} % <<< includes "figure ..."
%\usepackage{natbib}

\usepackage{subfig}

\providecommand{\btVFill}{\vskip0pt plus 1filll}

% \newcommand{\todo}[1]{} % uncomment this before publishing
\newcommand{\todo}[1]{\textcolor{red}{//\textbf{TODO: {#1}}} }

\newcommand{\eng}[1] {(angl. \textit{#1})}
\newcommand{\engp}[2] {(angl. \textit{#1}, toliau: #2)}

\usetikzlibrary{matrix,shapes,arrows,fit,tikzmark}
\newcolumntype{.}{D{.}{.}{-1}}

% Some options common to all the nodes and paths
\tikzset{   
        every picture/.style={remember picture,baseline},
        every node/.style={anchor=base,align=center,outer sep=1.5pt},
        every path/.style={thick},
        }

\newcommand\marktopleft[1]{%
    \tikz[overlay,remember picture] 
        \node (marker-#1-a) at (.1em,.3em) {};%
}
\newcommand\markbottomright[1]{%
    \tikz[overlay,remember picture] 
        \node (marker-#1-b) at (.1em,.3em) {};%
    \tikz[overlay,remember picture,inner sep=3pt]
        \node[draw=red,rounded corners,fit=(marker-#1-a.north west) (marker-#1-b.south east)] {};%
}

\newcolumntype{P}[1]{>{\centering\arraybackslash}p{#1}}

% TODO: attempt to remove the 1s
% TITLE PAGE 

\title[Privačios informacijos išsaugojimas]
{Privačios informacijos išsaugojimas taikant dirbtinio intelekto technologijas}
% \subtitle{Presentation}

\author[Paulius Milmantas]
{Paulius Milmantas\\
Darbo vadovas: dr. Linas Petkevičius}
 
\institute[VU-MIF] % (optional)
{
  Vilniaus Universitetas \\
  Matematikos ir informatikos fakultetas \\
}

\date[2021-05-14]{Bakalauro darbo gynimas}


% Let's get started
\begin{document}

\addtobeamertemplate{frametitle}{}{%
\begin{textblock*}{100mm}(.95\textwidth,-0.95cm)
\includegraphics[height=1cm,width=1cm]{VU_Logo_pilkas.png}
\end{textblock*}
}

%
\begin{frame}
  \titlepage
\end{frame}
%

% \begin{frame}[c]{Turinys}
%   \setcounter{tocdepth}{1}    
%   \tableofcontents
% \end{frame}


% Tiriama sritis
 \begin{frame}[c]{Tiriama sritis}
 
\setlength{\parindent}{10ex} Mašininis mokymas yra dirbtinio intelekto sritis, kuri pasitelkia statistinius algoritmus, kad apibrėžtų duomenų generavimo mechanizmą, ar egzistuojančius sąryšius, priklausomybes. 

 \end{frame}
% Problematika
\section{Problematika}

 \begin{frame}[c]{Problematika}

    \begin{enumerate}
        \item Turint sukurtą modelį, neturi būti galima atgaminti duomenų, pagal kuriuos jis buvo mokomas, bei negali būti identifikuoti asmenys. 
        \item Trečios šalys neturi matyti įvedamų duomenų. Tai gali būti tinklo saugumo spragos, duomenų surinkimo aplikacijų spragos ir t.t…
        \item Modelio išvesties neturi matyti asmenys, kuriems šie duomenys nepriklauso.
        \item Sukurtas modelis negali būti niekieno pasisavintas. 
    \end{enumerate}

 \end{frame}
% Modelių duomenų lyginimas (1)
\begin{frame}[c]{Modelių duomenų lyginimas (1)}
 
 \begin{equation}
     atvirumas(s[r])_{\theta} = log_{2}|r| - log_{2} rangas_{\theta}(s[r])
 \end{equation}
 
Naudojama teorijoje, dėl sunkiai apskaičiuojamo rango [2].

$\bf s$ - duomenų rinkinys. \\
$\bf r \in R$, parenkamas atsitiktinai.

\end{frame}
% Modelių duomenų lyginimas (2)
\begin{frame}[c]{Modelių duomenų lyginimas (2)}

\begin{equation}
    atvirumas(s[r])_{\theta} = -log_{2} \int_{0}^{Px_{\theta}(s[r])} \rho(x)dx
\end{equation}

Dėl grafinės interpretacijos naudojama praktikoje.

$\bf Px$ - logaritminis entropijos matas. \\
$\bf s[r]$ - entropija yra $\rho(.)$ pasiskirstymo distribucijos.

\end{frame}
% Tyrimo metodika
\begin{frame}[c]{Pasiūlyta tyrimo metodika}

\begin{equation}
    DMDK = {\sum_{n=0}^{m} ({\sum_{k=0}^{h} (max_{\epsilon}((|\epsilon| + D_{n, k}) : \epsilon \in R, modelis(|\epsilon| + D_{n, k}) = modelis(D_{n, k})))}/{h})}/{m}
\end{equation}

$\bf DMDK$ - Didžiausias modelio duomenų nuokrypis. \\
$\bf D_{eilut:n,stulp:k}$ - duomenys n eilutėje ir k stulpelyje. \\
$\bf \epsilon$ - ieškomas didžiausias galimas kintamasis, su kuriuo modelis nepakeičia išvesties rezultatų. \\
$\bf m$ - duomenų eilučių skaičius. \\
$\bf h$ - parametrų skaičius (stulpeliai).

\end{frame}

% \appendix

\end{document}
